\numoptstrue
\yearpaper{2018}{6 May 2018 \textperiodcentered{} Test Booklet Code AA \textperiodcentered{} Physics Q 1--45 \textperiodcentered{} 180 Marks}{}
\begin{multicols}{2}

\QQ{The volume ($V$) of a monatomic gas varies with its temperature ($T$), as shown in the graph. The ratio of work done by the gas, to the heat absorbed by it, when it undergoes a change from state A to state B, is
\fig{\gaxes{$T$}{$V$}{\draw (0,0) -- (1.9,0.9); \draw[very thick,->] (0.6,0.28) -- (1.1,0.52); \node[lbl,above] at (0.6,0.28) {A}; \node[lbl,above] at (1.4,0.66) {B}; \node[lbl,below left] at (0,0) {O};}}}
{\opt{$\dfrac13$}{$\dfrac23$}{$\dfrac25$}{$\dfrac27$}}

\QQ{The fundamental frequency in an open organ pipe is equal to the third harmonic of a closed organ pipe. If the length of the closed organ pipe is $20\unit{cm}$, the length of the open organ pipe is}
{\opt{12.5\,cm}{8\,cm}{13.2\,cm}{16\,cm}}

\QQ{At what temperature will the rms speed of oxygen molecules become just sufficient for escaping from the Earth's atmosphere? (Given: Mass of oxygen molecule $m = 2.76\times10^{-26}\unit{kg}$, Boltzmann's constant $k_B = 1.38\times10^{-23}\unit{J\,K^{-1}}$)}
{\opt{$5.016\times10^4\unit{K}$}{$8.360\times10^4\unit{K}$}{$2.508\times10^4\unit{K}$}{$1.254\times10^4\unit{K}$}}

\QQ{The efficiency of an ideal heat engine working between the freezing point and boiling point of water, is}
{\opt{6.25\%}{20\%}{26.8\%}{12.5\%}}

\QQ{A carbon resistor of $(47\pm4.7)\unit{k\Omega}$ is to be marked with rings of different colours for its identification. The colour code sequence will be}
{\opt{Yellow -- Green -- Violet -- Gold}{Yellow -- Violet -- Orange -- Silver}{Violet -- Yellow -- Orange -- Silver}{Green -- Orange -- Violet -- Gold}}

\QQ{A set of `$n$' equal resistors, of value `$R$' each, are connected in series to a battery of emf `$E$' and internal resistance `$R$'. The current drawn is $I$. Now, the `$n$' resistors are connected in parallel to the same battery. Then the current drawn from battery becomes $10I$. The value of `$n$' is}
{\opt{20}{11}{10}{9}}

\QQ{A battery consists of a variable number `$n$' of identical cells (having internal resistance `$r$' each) which are connected in series. The terminals of the battery are short-circuited and the current $I$ is measured. Which of the graphs shows the correct relationship between $I$ and $n$?}
{\optfig{\gaxes{$n$}{$I$}{\draw (0,0) .. controls (0.2,0.9) and (0.6,1.3) .. (1.3,1.4);}}
{\gaxes{$n$}{$I$}{\draw (0,0) -- (1.5,1.3);}}
{\gaxes{$n$}{$I$}{\draw (0,1.2) -- (1.4,1.2);}}
{\gaxes{$n$}{$I$}{\draw (0,0.1) .. controls (0.9,0.15) and (1.3,0.6) .. (1.4,1.5);}}}

\QQ{Unpolarised light is incident from air on a plane surface of a material of refractive index `$\mu$'. At a particular angle of incidence `$i$', it is found that the reflected and refracted rays are perpendicular to each other. Which of the following options is correct for this situation?}
{\opt{$i = \sin^{-1}\!\left(\dfrac1\mu\right)$}{Reflected light is polarised with its electric vector perpendicular to the plane of incidence}{Reflected light is polarised with its electric vector parallel to the plane of incidence}{$i = \tan^{-1}\!\left(\dfrac1\mu\right)$}}

\QQ{In Young's double slit experiment the separation $d$ between the slits is $2\unit{mm}$, the wavelength $\lambda$ of the light used is $5896\unit{\text{\AA}}$ and distance $D$ between the screen and slits is $100\unit{cm}$. It is found that the angular width of the fringes is $0.20^\circ$. To increase the fringe angular width to $0.21^\circ$ (with same $\lambda$ and $D$) the separation between the slits needs to be changed to}
{\opt{2.1\,mm}{1.9\,mm}{1.8\,mm}{1.7\,mm}}

\QQ{An astronomical refracting telescope will have large angular magnification and high angular resolution, when it has an objective lens of}
{\opt{large focal length and large diameter}{large focal length and small diameter}{small focal length and large diameter}{small focal length and small diameter}}

\QQ{The ratio of kinetic energy to the total energy of an electron in a Bohr orbit of the hydrogen atom, is}
{\opt{$2:-1$}{$1:-1$}{1:1}{$1:-2$}}

\QQ{An electron of mass $m$ with an initial velocity $\vec V = V_0\hat i$ ($V_0>0$) enters an electric field $\vec E = -E_0\hat i$ ($E_0$ = constant $>0$) at $t = 0$. If $\lambda_0$ is its de-Broglie wavelength initially, then its de-Broglie wavelength at time $t$ is}
{\opt{$\lambda_0t$}{$\lambda_0\left(1+\dfrac{eE_0}{mV_0}t\right)$}{$\dfrac{\lambda_0}{\left(1+\dfrac{eE_0}{mV_0}t\right)}$}{$\lambda_0$}}

\QQ{For a radioactive material, half-life is 10 minutes. If initially there are 600 number of nuclei, the time taken (in minutes) for the disintegration of 450 nuclei is}
{\opt{30}{10}{20}{15}}

\QQ{When the light of frequency $2\nu_0$ (where $\nu_0$ is threshold frequency), is incident on a metal plate, the maximum velocity of electrons emitted is $v_1$. When the frequency of the incident radiation is increased to $5\nu_0$, the maximum velocity of electrons emitted from the same plate is $v_2$. The ratio of $v_1$ to $v_2$ is}
{\opt{4:1}{1:4}{1:2}{2:1}}

\QQ{In the circuit shown in the figure, the input voltage $V_i$ is 20\,V, $V_{BE} = 0$ and $V_{CE} = 0$. The values of $I_B$, $I_C$ and $\beta$ are given by
\fig{\begin{circuitikz}[scale=0.62,transform shape,font=\Large]
\draw (3,0) node[npn](Q){}; \draw (Q.B) to[R,l_={$R_B$ 500\,k$\Omega$}] ++(-2.4,0) node[circ]{} node[left]{$V_i$};
\draw (Q.C) to[R,l_={$R_C$ 4\,k$\Omega$}] ++(0,2) node[circ]{} node[above]{20\,V}; \draw (Q.E) -- ++(0,-0.4) node[ground]{};
\end{circuitikz}}}
{\opt{$I_B = 20\unit{\mu A}$, $I_C = 5\unit{mA}$, $\beta = 250$}{$I_B = 25\unit{\mu A}$, $I_C = 5\unit{mA}$, $\beta = 200$}{$I_B = 40\unit{\mu A}$, $I_C = 10\unit{mA}$, $\beta = 250$}{$I_B = 40\unit{\mu A}$, $I_C = 5\unit{mA}$, $\beta = 125$}}

\QQ{In a p-n junction diode, change in temperature due to heating}
{\opt{does not affect resistance of p-n junction}{affects only forward resistance}{affects only reverse resistance}{affects the overall V--I characteristics of p-n junction}}

\QQ{In the combination of the following gates the output Y can be written in terms of inputs A and B as
\fig{\begin{circuitikz}[scale=0.7,transform shape]
\draw (3,1.2) node[and port](G1){}; \draw (3,-1.0) node[and port](G2){}; \draw (5.2,0.1) node[or port](G3){};
\draw (0,1.6) node[circ]{} node[left]{A} -- (0.6,1.6) coordinate(a1); \draw (a1) |- (G1.in 1);
\draw (a1) |- (1.0,-0.8) node[not port,anchor=in,scale=0.6](N2){}; \draw (N2.out) -| ($(G2.in 1)+(-0.2,0)$) -- (G2.in 1);
\draw (0,-1.6) node[circ]{} node[left]{B} -- (0.3,-1.6) coordinate(b1); \draw (b1) |- (G2.in 2);
\draw (b1) |- (1.0,1.0) node[not port,anchor=in,scale=0.6](N1){}; \draw (N1.out) -| ($(G1.in 2)+(-0.2,0)$) -- (G1.in 2);
\draw (G1.out) -| (G3.in 1); \draw (G2.out) -| (G3.in 2); \draw (G3.out) -- ++(0.4,0) node[circ]{} node[right]{Y};
\end{circuitikz}}}
{\opt{$\overline{A\cdot B}+A\cdot B$}{$A\cdot\bar B+\bar A\cdot B$}{$\overline{A\cdot B}$}{$\overline{A+B}$}}

\QQ{An em wave is propagating in a medium with a velocity $\vec V = V\hat i$. The instantaneous oscillating electric field of this em wave is along $+y$ axis. Then the direction of oscillating magnetic field of the em wave will be along}
{\opt{$-y$ direction}{$+z$ direction}{$-z$ direction}{$-x$ direction}}

\QQ{The refractive index of the material of a prism is $\sqrt2$ and the angle of the prism is $30^\circ$. One of the two refracting surfaces of the prism is made a mirror inwards, by silver coating. A beam of monochromatic light entering the prism from the other face will retrace its path (after reflection from the silvered surface) if its angle of incidence on the prism is}
{\opt{$30^\circ$}{$45^\circ$}{$60^\circ$}{zero}}

\QQ{An object is placed at a distance of $40\unit{cm}$ from a concave mirror of focal length $15\unit{cm}$. If the object is displaced through a distance of $20\unit{cm}$ towards the mirror, the displacement of the image will be}
{\opt{30\,cm towards the mirror}{36\,cm away from the mirror}{30\,cm away from the mirror}{36\,cm towards the mirror}}

\QQ{The magnetic potential energy stored in a certain inductor is $25\unit{mJ}$, when the current in the inductor is $60\unit{mA}$. This inductor is of inductance}
{\opt{1.389\,H}{138.88\,H}{0.138\,H}{13.89\,H}}

\QQ{An electron falls from rest through a vertical distance $h$ in a uniform and vertically upward directed electric field E. The direction of electric field is now reversed, keeping its magnitude the same. A proton is allowed to fall from rest in it through the same vertical distance $h$. The time of fall of the electron, in comparison to the time of fall of the proton is}
{\opt{10 times greater}{5 times greater}{smaller}{equal}}

\QQ{The electrostatic force between the metal plates of an isolated parallel plate capacitor C having a charge Q and area A, is}
{\opt{proportional to the square root of the distance between the plates.}{linearly proportional to the distance between the plates.}{independent of the distance between the plates.}{inversely proportional to the distance between the plates.}}

\QQ{A tuning fork is used to produce resonance in a glass tube. The length of the air column in this tube can be adjusted by a variable piston. At room temperature of $27^\circ$C two successive resonances are produced at $20\unit{cm}$ and $73\unit{cm}$ of column length. If the frequency of the tuning fork is 320\,Hz, the velocity of sound in air at $27^\circ$C is}
{\opt{350\,m/s}{339\,m/s}{330\,m/s}{300\,m/s}}

\QQ{A pendulum is hung from the roof of a sufficiently high building and is moving freely to and fro like a simple harmonic oscillator. The acceleration of the bob of the pendulum is $20\unit{m/s^2}$ at a distance of $5\unit{m}$ from the mean position. The time period of oscillation is}
{\opt{2\,s}{$\pi$\,s}{$2\pi$\,s}{1\,s}}

\QQ{A metallic rod of mass per unit length $0.5\unit{kg\,m^{-1}}$ is lying horizontally on a smooth inclined plane which makes an angle of $30^\circ$ with the horizontal. The rod is not allowed to slide down by flowing a current through it when a magnetic field of induction $0.25\unit{T}$ is acting on it in the vertical direction. The current flowing in the rod to keep it stationary is}
{\opt{14.76\,A}{5.98\,A}{7.14\,A}{11.32\,A}}

\QQ{A thin diamagnetic rod is placed vertically between the poles of an electromagnet. When the current in the electromagnet is switched on, then the diamagnetic rod is pushed up, out of the horizontal magnetic field. Hence the rod gains gravitational potential energy. The work required to do this comes from}
{\opt{the lattice structure of the material of the rod}{the magnetic field}{the current source}{the induced electric field due to the changing magnetic field}}

\QQ{An inductor $20\unit{mH}$, a capacitor $100\unit{\mu F}$ and a resistor $50\unit{\Omega}$ are connected in series across a source of emf, $V = 10\sin314t$. The power loss in the circuit is}
{\opt{2.74\,W}{0.43\,W}{0.79\,W}{1.13\,W}}

\QQ{Current sensitivity of a moving coil galvanometer is 5\,div/mA and its voltage sensitivity (angular deflection per unit voltage applied) is 20\,div/V. The resistance of the galvanometer is}
{\opt{250\,$\Omega$}{25\,$\Omega$}{40\,$\Omega$}{500\,$\Omega$}}

\QQ{A body initially at rest and sliding along a frictionless track from a height $h$ (as shown in the figure) just completes a vertical circle of diameter $AB = D$. The height $h$ is equal to
\fig{\begin{tikzpicture}[scale=0.7]
\draw (0,0) -- (3,0); \draw (0,1.8) -- (3,0); \fill (0,1.8) circle (1.5pt); \draw[<->] (-0.2,0) -- (-0.2,1.8) node[midway,left,lbl]{$h$};
\draw (3,0.4) circle (0.4); \node[lbl,below] at (3,0) {A}; \node[lbl,above] at (3,0.8) {B}; \draw[densely dashed] (3,0) -- (3,0.8);
\end{tikzpicture}}}
{\opt{$\dfrac75D$}{$D$}{$\dfrac32D$}{$\dfrac54D$}}

\QQ{Three objects, A: (a solid sphere), B: (a thin circular disk) and C: (a circular ring), each have the same mass M and radius R. They all spin with the same angular speed $\omega$ about their own symmetry axes. The amounts of work ($W$) required to bring them to rest, would satisfy the relation}
{\opt{$W_B>W_A>W_C$}{$W_A>W_B>W_C$}{$W_C>W_B>W_A$}{$W_A>W_C>W_B$}}

\QQ{A moving block having mass $m$, collides with another stationary block having mass $4m$. The lighter block comes to rest after collision. When the initial velocity of the lighter block is $v$, then the value of coefficient of restitution ($e$) will be}
{\opt{0.8}{0.25}{0.5}{0.4}}

\QQ{Which one of the following statements is \textit{incorrect}?}
{\opt{Frictional force opposes the relative motion.}{Limiting value of static friction is directly proportional to normal reaction.}{Rolling friction is smaller than sliding friction.}{Coefficient of sliding friction has dimensions of length.}}

\QQ{A toy car with charge $q$ moves on a frictionless horizontal plane surface under the influence of a uniform electric field $\vec E$. Due to the force $q\vec E$, its velocity increases from 0 to $6\unit{m/s}$ in one second duration. At that instant the direction of the field is reversed. The car continues to move for two more seconds under the influence of this field. The average velocity and the average speed of the toy car between 0 to 3 seconds are respectively}
{\opt{1\,m/s, 3.5\,m/s}{1\,m/s, 3\,m/s}{2\,m/s, 4\,m/s}{1.5\,m/s, 3\,m/s}}

\QQ{A block of mass $m$ is placed on a smooth inclined wedge ABC of inclination $\theta$ as shown in the figure. The wedge is given an acceleration `$a$' towards the right. The relation between $a$ and $\theta$ for the block to remain stationary on the wedge is
\fig{\begin{tikzpicture}[scale=0.8]
\fill[pattern=north east lines] (0,0) -- (0,1.5) -- (2.4,0) -- cycle; \draw (0,0) -- (0,1.5) -- (2.4,0) -- cycle; \draw (-0.4,0) -- (2.8,0);
\node[lbl,above left] at (0,1.5) {A}; \node[lbl,below] at (0,0) {C}; \node[lbl,below] at (2.4,0) {B};
\draw[rotate around={-32:(1.1,0.81)}] (0.85,0.81) rectangle (1.35,1.15); \node[lbl] at (1.45,1.25) {$m$};
\draw (1.9,0) arc (180:148:0.5); \node[lbl] at (1.75,0.18) {$\theta$}; \draw[->] (2.5,0.6) -- (3.2,0.6) node[right,lbl]{$a$};
\end{tikzpicture}}}
{\opt{$a = g\cos\theta$}{$a = \dfrac g{\sin\theta}$}{$a = \dfrac g{\operatorname{cosec}\theta}$}{$a = g\tan\theta$}}

\QQ{The moment of the force, $\vec F = 4\hat i+5\hat j-6\hat k$ at $(2,0,-3)$, about the point $(2,-2,-2)$, is given by}
{\opt{$-7\hat i-8\hat j-4\hat k$}{$-4\hat i-\hat j-8\hat k$}{$-8\hat i-4\hat j-7\hat k$}{$-7\hat i-4\hat j-8\hat k$}}

\QQ{A student measured the diameter of a small steel ball using a screw gauge of least count $0.001\unit{cm}$. The main scale reading is $5\unit{mm}$ and zero of circular scale division coincides with 25 divisions above the reference level. If screw gauge has a zero error of $-0.004\unit{cm}$, the correct diameter of the ball is}
{\opt{0.053\,cm}{0.525\,cm}{0.521\,cm}{0.529\,cm}}

\QQ{A solid sphere is rotating freely about its symmetry axis in free space. The radius of the sphere is increased keeping its mass same. Which of the following physical quantities would remain constant for the sphere?}
{\opt{Rotational kinetic energy}{Moment of inertia}{Angular velocity}{Angular momentum}}

\QQ{The kinetic energies of a planet in an elliptical orbit about the Sun, at positions A, B and C are $K_A$, $K_B$ and $K_C$, respectively. AC is the major axis and SB is perpendicular to AC at the position of the Sun S as shown in the figure. Then
\fig{\begin{tikzpicture}[scale=0.7]
\draw (0,0) ellipse (2 and 0.9); \draw (-2,0) -- (2,0); \fill (-1.2,0) circle (1.5pt) node[below right,lbl]{S}; \draw (-1.2,0) -- (-1.2,0.72);
\node[lbl,left] at (-2,0) {A}; \node[lbl,right] at (2,0) {C}; \node[lbl,above] at (-1.2,0.72) {B};
\end{tikzpicture}}}
{\opt{$K_B<K_A<K_C$}{$K_A>K_B>K_C$}{$K_A<K_B<K_C$}{$K_B>K_A>K_C$}}

\QQ{If the mass of the Sun were ten times smaller and the universal gravitational constant were ten times larger in magnitude, which of the following is \textit{not} correct?}
{\opt{Time period of a simple pendulum on the Earth would decrease.}{Walking on the ground would become more difficult.}{Raindrops will fall faster.}{`$g$' on the Earth will not change.}}

\QQ{A solid sphere is in rolling motion. In rolling motion a body possesses translational kinetic energy ($K_t$) as well as rotational kinetic energy ($K_r$) simultaneously. The ratio $K_t:(K_t+K_r)$ for the sphere is}
{\opt{10:7}{5:7}{7:10}{2:5}}

\QQ{A small sphere of radius `$r$' falls from rest in a viscous liquid. As a result, heat is produced due to viscous force. The rate of production of heat when the sphere attains its terminal velocity, is proportional to}
{\opt{$r^5$}{$r^2$}{$r^3$}{$r^4$}}

\QQ{The power radiated by a black body is P and it radiates maximum energy at wavelength, $\lambda_0$. If the temperature of the black body is now changed so that it radiates maximum energy at wavelength $\dfrac34\lambda_0$, the power radiated by it becomes $nP$. The value of $n$ is}
{\opt{$\dfrac{256}{81}$}{$\dfrac43$}{$\dfrac34$}{$\dfrac{81}{256}$}}

\QQ{Two wires are made of the same material and have the same volume. The first wire has cross-sectional area A and the second wire has cross-sectional area 3A. If the length of the first wire is increased by $\Delta l$ on applying a force F, how much force is needed to stretch the second wire by the same amount?}
{\opt{4F}{6F}{9F}{F}}

\QQ{A sample of 0.1\,g of water at $100^\circ$C and normal pressure ($1.013\times10^5\unit{N\,m^{-2}}$) requires 54\,cal of heat energy to convert to steam at $100^\circ$C. If the volume of the steam produced is 167.1\,cc, the change in internal energy of the sample, is}
{\opt{42.2\,J}{208.7\,J}{104.3\,J}{84.5\,J}}

\end{multicols}
